% Lösungen Einheit 3 von 5 – Normale (zusammenbau v0.1)
\einheitenkopf{Einheit 3 von 5 · Normale}

\erg{18}{a) Normale, Steigung $-\frac{1}{4}$ – senkrecht heißt negativer Kehrwert. \quad b) $m_n = -\frac{1}{2} = -0{,}5$; $n(x) = -0{,}5 \cdot (x - 3) + 1 = -0{,}5x + 2{,}5$ \quad c) $y_P = f(1) = 3$; Anstieg $f'(1) = 3 - 2 + 1 = 2$; $m_n = -\frac{1}{2}$ (nicht $2$); $n(x) = -0{,}5 \cdot (x - 1) + 3 = -0{,}5x + 3{,}5$ \quad d) $m = -\frac{1}{2}$; $y = -0{,}5 \cdot (x - 4) + 1 = -0{,}5x + 3$ \quad e) $f(4) = 3$, $f'(4) = 1$; Normale $n(x) = -x + 7$, Nullstelle $x = 7$; Länge $\sqrt{3^2 + 3^2} = \sqrt{18} \approx 4{,}24$ m (nicht $f(4) = 3$) \quad f) Lot $y = -\frac{1}{3}(x - 10)$; Schnitt mit $y = 3x$: $x = 1$, Lotfußpunkt $F(1 | 3)$; Abstand $\sqrt{9^2 + 3^2} = \sqrt{90} \approx 9{,}49$ \quad g) Steigung des Graphen von $f'$ in $A$: $f''(2) = 2$ (nicht $f'(2)$); Normale $y = -\frac{1}{2}(x - 2) + 2$; Schnitt mit der $y$-Achse: $y = 3$, Mittelpunkt $(0 | 3)$ \quad h) Der kürzeste Abstand des Punktes $(3 | 2)$ vom Graphen von $h$: der Bruch ist die Steigung der Verbindung von $(3 | 2)$ zum Graphenpunkt $(x | h(x))$; Produkt $-1$ heißt, sie steht senkrecht auf der Tangente, liegt also auf der Normalen – der Graphenpunkt ist der Lotfußpunkt, der Abstand die Länge dieser Verbindung (nicht $h(3)$).}
\erg{19}{a) Mira dreht nur das Vorzeichen. Richtig: $m_n = -\frac{1}{3}$; $n(x) = -\frac{1}{3}(x - 3) + 2 = -\frac{1}{3}x + 3$. \quad b) Jede andere Verbindung zur Geraden ist Hypotenuse eines rechtwinkligen Dreiecks, dessen eine Kathete das Lot ist – also länger als das Lot.}
